\hypertarget{the-command-surface}{%
\chapter{5. The command surface}\label{the-command-surface}}

Chapter 4 drew the line around the trusted base and noted how much sits
outside it: the operator CLI, the compiler, the argument parser, the
terminal filter. The command surface is that outside, the operator's
whole interface to the system, and the fact that nearly all of it falls
beyond the trusted line is the shape worth seeing. The entry point holds
no authority, the wire to the daemon is small enough to read in a
sitting, and the work that looks heavy, compiling and signing a policy,
happens here, on the untrusted side, where a bug is the operator's
problem and not the confinement's.

\hypertarget{the-shim}{%
\section{5.1 The shim}\label{the-shim}}

The command a user or an agent types is \texttt{kennel}, and it is a
13-line static binary that holds no authority and does no work of its
own. It detects its context and execs the right execution unit,
replacing itself: the host execution unit on the host, the in-cage
spawn-requester inside a kennel. What makes it small is that the
detection is not a check. The shim execs the host unit first,
unconditionally; on the host that succeeds and the host unit takes over,
and inside a kennel it fails, because the host unit is absent from the
kennel's view and not in its \texttt{exec.allow}, so the shim falls
through to the spawn unit. The failure to reach the host unit is the
context signal (\texttt{construction-by-absence}).

That the signal is a failure rather than a probe is what keeps the shim
out of the security argument entirely. A wrong guess cannot reach host
authority, because the thing a wrong guess would reach, the host unit,
is unreachable from inside a kennel by construction, allowed nowhere and
soon absent from every cage's view. The dispatch is therefore ergonomics
and not a boundary: it decides which of two units runs, and the one that
must not run from inside cannot be run from inside no matter what the
shim does (\texttt{do-less}). One static artifact serves both sides, so
there is nothing to install differently in a kennel than on the host.

\hypertarget{outside-the-trusted-base}{%
\section{5.2 Outside the trusted base}\label{outside-the-trusted-base}}

The host execution unit is the operator CLI, its own crate, outside the
daemon's trusted base. It links the control-wire types and nothing of
the daemon's enforcement code, so the dependencies that come with an
operator tool, an argument parser, a JSON reader for the trust manifest,
the terminal-escape filter, stay out of the daemon's closure entirely.
The protocol cannot drift between the two, because both the CLI and the
daemon depend on the one control crate for the wire types: one
definition, linked twice (\texttt{rule-of-1}).

The CLI is stateless and splits cleanly in two. Some verbs contact the
daemon over its control socket, the ones that start, attach to, stop, or
list a running kennel. The rest operate only on local files and never
open the socket at all: the whole policy group, the audit query, key
generation, the trust-manifest review. This is the runtime-and-compiler
split of chapter 4 seen from the operator's side. The compiler is the
largest single thing the CLI carries, and it runs here, locally, turning
a source policy into the signed settled artefact, while the daemon links
none of it and only ever verifies the artefact's signature and acts on
it. The heavy authoring work and the light enforcement work are on
opposite sides of the trusted line, and the command surface is where the
heavy side lives.

\hypertarget{the-control-protocol}{%
\section{5.3 The control protocol}\label{the-control-protocol}}

The wire to the daemon is deliberately small. The socket is the user's
own, mode 0600 under the runtime directory, and the first thing across
it is a version handshake: the client sends its schema version and
build, and the daemon answers compatible or, on a mismatch, drops the
connection without reading a request. A schema skew is refused before
any request body or policy is parsed, so a client and a daemon that
disagree about the protocol never get far enough to misread each other
(\texttt{refuse-to-start}).

The request and reply bodies are fixed-layout struct messages, each a
leading op or tag byte followed by primitively encoded fields. There is
no JSON, no request identifier, and no method-name string: the op byte
is the method selector. Integers are native-endian, strings and lists
are length-prefixed and bounded, a string at 64 kilobytes and a list at
4,096 elements, and UTF-8 is validated on decode. The verbs are few,
start a kennel, stop one, list them, and answer the SSH bastion's
authorized-keys query, and their replies as few. An error is a
first-class reply carrying a message, not a coded catalogue to keep in
sync; a malformed frame, an unknown op or a truncated, oversized, or
non-UTF-8 field, is refused at decode and closes the connection. A wire
this narrow, with every field bounded and the selector a single byte, is
one whose decoder makes most malformed input unrepresentable rather than
handled (\texttt{make-invalid-unrepresentable}).

What the protocol does not carry is the workload's bytes. The standard
streams travel as descriptors, passed alongside the start frame as
ancillary data, not copied through the daemon: a non-interactive run
hands over its three stdio descriptors and the workload is attached to
them directly. An interactive run hands over a single socket, and the
spawn seal allocates a controlling pty inside the kennel's own devpts,
so the workload has real job control and the operator's own terminal is
never exposed to it. The daemon's broker holds the pty master and
proxies it, the CLI raw-modes the terminal and proxies bytes both ways,
and the filtering of dangerous terminal escapes in the workload's output
runs in the CLI, on the bytes' way to the real terminal, so the parser
of workload-controlled output never executes in the daemon
(\texttt{control-not-data-plane}). The daemon decides and routes; the
bytes flow past it.

\hypertarget{the-verbs}{%
\section{5.4 The verbs}\label{the-verbs}}

The surface above the wire is a small set of verbs dispatched on the
first argument, each parsing its own flags, with the command list
rendered from an in-binary table so the help cannot drift from what
exists. The lifecycle verbs are the ones that reach the daemon:
\texttt{run} starts a kennel from a policy and runs a workload in it,
foreground and, when interactive, detachable; \texttt{attach} reconnects
a terminal to a running one; \texttt{stop} and \texttt{list} do as they
say. A \texttt{run} from a pre-compiled settled policy needs no key,
because the daemon verifies the artefact's signature against its trust
store; only a run straight from source compiles and signs in memory and
so needs a signing key.

Everything else is local. The policy group is the authoring surface,
compiling, validating, signing, linting, diffing, evaluating a policy
against the threat catalogue, and re-pinning a template to a newer
version, all on files, none touching the daemon. \texttt{keygen} makes a
signing key, \texttt{audit} queries a kennel's recorded account, and
\texttt{review} re-pins a workspace's trust manifest, the deliberate
operator act the run path will not perform as a side effect (T2.8). The
weight of the command surface is here, in authoring and inspection, on
the untrusted side of the line, which is the point: the operator's tool
is rich and the trusted base it drives is not.
